\chapter*{Abstract}
\label{ch:abstract}

The growing generation of \glsauto{fw} and the surplus of \glsauto{cgl} from biodiesel production represent two significant organic waste streams requiring sustainable management solutions. This study evaluated the technical feasibility of a single-stage anaerobic co-digestion system treating \glsauto{fw} and \glsauto{cgl} as co-substrates, with the objective of verifying biogas production while providing a valorisation route for \glsauto{cgl}. The investigation was conducted in a \qty{1}{\litre} \glsauto{cstr} operated under mesophilic conditions (\qty{35}{\celsius}) with a \glsauto{hrt} of \num{30} days. \Glsentrylongauto{fw} was collected from a university canteen, and \glsauto{cgl} was obtained from biodiesel production residues; sludge from a anaerobic reactor from a municipal \glsauto{wwtp} served as inoculum. The reactor was operated in sequential phases: 1-an initial adaptation period, 2-co-digestion of \glsauto{fw} and \glsauto{cgl}, 3-\glsauto{cgl} only feeding, and 4-a post-feding period. Biogas production was monitored continuously using \glsauto{ampts}, while process stability was assessed through measurements of pH, total \glsauto{alk}, \glsauto{vfa}, and their \glsauto{vfa_alk}. Substrate characterisation included \glsauto{ts}, \glsauto{vs}, density, \glsauto{tc}, and \glsauto{tn} analyses. Results demonstrated that Phase 2 (co-digestion of \glsauto{fw} and \glsauto{cgl}) yielded the highest \glsauto{ch4} production rates and cumulative volumes, with flow spikes immediately following feeding events indicating rapid substrate conversion. In contrast, pahse 3 (\glsauto{cgl}-only feeding) resulted in reduced biogas production despite continued methanogenic activity. Total \glsauto{alk} remained elevated following bicarbonate supplementation, maintaining pH around \num{7} throughout most of the experiment; \glsauto{vfa} concentrations fluctuated between \numrange{3000}{5000} \unit{\milli\gram\CHCOOH\per\litre}, with \glsauto{vfa_alk} of \numrange{0.3}{1.2}. These findings support the hypothesis that \glsauto{cgl} can be useful for biogas production when co-digested with \glsauto{fw} or alone. The study demonstrates the technical viability of anaerobic co-digestion as a strategy for valorising \glsauto{cgl} within a circular bioeconomy framework, though further research is needed to optimise \glsauto{cgl} loading rates and develop robust feeding strategies for industrial application.

\vspace{1em}
\noindent\textbf{Keywords:} \Keywords.
